\section{Finite-width representation learning}
\label{sec:dntk}

Infinite-width NTK dynamics are linear in function space: $f(t)$ stays in the affine space $f(0)+\mathrm{span}\{K(x_\mu,\,\cdot\,)\}$. The first correction that moves the kernel is $O(L/N)$ at criticality~\cite{roberts2022pdlt} and $O(\gamma_0)$ at infinite width in $\mu$P.

\subsection{Quadratic model}

Taylor-expand in $\delta\theta=\theta-\theta_0$:
\begin{equation}
  f(x;\theta_0+\delta\theta)
  =f(x;\theta_0)
  +J(x)\cdot\delta\theta
  +\tfrac12\,\delta\theta^\top H(x)\,\delta\theta
  +O(\|\delta\theta\|^3),
  \label{eq:taylor}
\end{equation}
with $J(x)=\partial_\theta f(x)|_{\theta_0}$ and $H_{pq}(x)=\partial_{\theta_p}\partial_{\theta_q}f(x)|_{\theta_0}$. The first two terms are a linear model with features $J$; their Gram matrix is the NTK
\begin{equation}
  \Theta(x,x')=J(x)\cdot J(x').
  \label{eq:ntk-again}
\end{equation}
Truncating at linear order in $\delta\theta$ recovers frozen-NTK training. The quadratic term is the leading representation-learning correction.

\subsection{Differential of the NTK}

Along a parameter direction $u$,
\begin{equation}
  \bigl(\nabla_u\Theta\bigr)(x,x')
  =H(x)\,u\cdot J(x')
  +J(x)\cdot H(x')\,u.
  \label{eq:dntk}
\end{equation}
The trilinear kernel
\begin{equation}
  D(x,x',x'')
  \equiv
  \partial_\theta\Theta(x,x')\cdot\partial_\theta f(x'')
  \label{eq:dntk-def}
\end{equation}
is the dNTK~\cite{roberts2022pdlt}. It vanishes for a linear model ($H=0$). In NTK scaling $\delta\theta=O(N^{-1/2})$, so $D$ does not move $\Theta$ by an $O(1)$ amount. At finite $N$, one GD step of size $\eta$ updates the kernel by
\begin{equation}
  \Delta\Theta(x,x')
  =-\frac{\eta}{P}\sum_{\mu} D(x,x',x_\mu)\,\Delta_\mu
  +O(\eta^2).
  \label{eq:dtheta}
\end{equation}

The first layer of an MLP contributes $D^{(1)}=0$ (the map $x\mapsto W^{(1)}x$ is linear in $W^{(1)}$). Each subsequent layer sources a nonzero $D$ from $\ddot\ph$ and from products of earlier $\Theta$'s. After $L$ layers at criticality, $D=O(L)$, hence
\begin{equation}
  \frac{\Delta\Theta}{\Theta}
  =O\!\left(\frac{L}{N}\right)
  \label{eq:dtheta-LN}
\end{equation}
for NTK-scaled GD with $\eta=O(1)$. The same ratio is $O(1)$ in $\mu$P at infinite $N$.

\subsection{Nearly-kernel prediction}

Keeping the quadratic term in~\eqref{eq:taylor} and going to a stationary point of square loss, the trained predictor is, to first order in $D$,
\begin{equation}
  f(x;\infty)
  =f_{\mathrm{NTK}}(x)
  +\frac1{2P^2}\sum_{\mu\nu}
  D(x,x_\mu,x_\nu)\,c_{\mu\nu}
  +O\bigl((L/N)^2\bigr),
  \label{eq:nearly-kernel}
\end{equation}
where $c_{\mu\nu}$ is bilinear in the training residuals of the linear problem~\cite{roberts2022pdlt}. Predictions are still determined by a finite collection of kernels $(\Theta,D)$.

The kernel picture is accurate when $L/N\ll 1$ and $d<\tfrac12$. Representation learning is $O(1)$ when $L/N=O(1)$, or when $d=\tfrac12$ with $\gamma_0=O(1)$. At fixed compute the useful aspect ratio is an $O(1)$ value of $L/N$: large enough that $D$ can move features, small enough that the $L/N$ expansion remains under control.

\begin{exercise}
For $f(x)=\tfrac12 A x^2$ with scalar $x$ and parameter $A$, compute $\Theta$ and $D$, and verify~\eqref{eq:dtheta} for one GD step.
\end{exercise}
